\section{本章小结}
\label{sec:nlipm-summary}

本章开头的问题是，约束变弯之后，线性规划的路径跟踪还剩下什么。骨架都留下了：在不等式上加对数障碍，驻定条件给出扰动互补$s_i\lambda_i=\mu$，让$\mu$减小到零。等式约束没有内部，仍要显式保留。

新增的内容来自曲率。障碍的Hessian多了$\nabla^2g_i/(-g_i)$一项，Newton系统左上角多了Lagrange Hessian，松弛方程的残差不再按$1-\alpha$的比例下降。非凸时，矩阵可能在切空间上不正定，解方程得到的方向可能是上升方向，要靠惯性检查和矩阵正则化修正，再由残差检验或过滤器决定接受哪一步。

内部性、中心性和最优性是三种不同的检查对象。$\bm{s},\bm{\lambda}>\bm{0}$由到边界的比例规则维护，中心性由$\sigma$和邻域控制，最优性只能由驻定残差、约束残差和互补量共同判断。很小的$\mu$或很小的$E_\mu$都不是原问题的结论，最后要按未经修改的问题报告。

起点有两种。障碍路径需要严格可行点，可以用辅助问题求出；原始—对偶迭代只要求$\bm{s},\bm{\lambda}>\bm{0}$，可以从非零残差出发。没有严格内点、约束重复或线性相关时，问题出在模型上，需要单独检查。

二阶锥和半正定锥有各自的障碍$-\log(t^2-\lVert\bm{u}\rVert_2^2)$和$-\log\det\bm{X}$，路径和间隙的形式不变，$n$换成障碍参数$\vartheta$。步长由二次方程或特征值给出。矩阵互补不能照搬标量操作，$\Delta\bm{X}$要经缩放保持对称，内积配对也要前后一致。

能给出什么结论，由模型的条件决定。凸问题上有整条路径、间隙$m\mu$、近似点的误差界，自协调时还有$\sqrt{\vartheta}$型的迭代界。一般非凸问题只有KKT点附近的局部路径和局部收敛，路径可以断开，也可能通向较差的解，线性规划的多项式结论不能外推。

计算上，每轮的主要代价是分解Jacobian和Lagrange Hessian组成的鞍点矩阵。压缩形式维数小但可能稠密，鞍点形式保留稀疏性，两者的条件数都按$1/\mu$增长。
